\chapter{Dataset Description}
\label{chap:dataset}

\section{Source and Collection}

The dataset for Type 1 Diabetes comes from a systematic immune profiling study of T1D PBMCs. Peripheral blood was collected from:
\begin{itemize}
    \item \textbf{46 T1D patients} --- confirmed diagnosis based on autoantibody positivity and insulin dependence
    \item \textbf{31 healthy controls} --- age/sex-matched, no autoimmune disease history
\end{itemize}

PBMCs were isolated by Ficoll density centrifugation and immediately processed for 10x Genomics Chromium single-cell RNA-seq (v3 chemistry). Sequencing was performed to a depth of $\sim$20,000 reads per cell.

\section{Data Specifications}

Table~\ref{tab:dataset_t1d} summarises the key properties of the dataset.

\begin{table}[h!]
\centering
\caption{T1D scRNA-seq Dataset Properties}
\label{tab:dataset_t1d}
\begin{tabular}{ll}
\toprule
\textbf{Property} & \textbf{Value} \\
\midrule
Total cells (post-QC) & 117,737 \\
Genes measured & 20,667 \\
Total patients/samples & 77 \\
T1D patients & 46 \\
Healthy controls & 31 \\
Cells per patient (mean) & 1,529 \\
Normalisation applied & Seurat SCTransform (SCT) \\
Matrix format & Sparse MTX (10x Genomics compatible) \\
\bottomrule
\end{tabular}
\end{table}

\section{Clinical Metadata}

The dataset includes rich clinical metadata for each patient (Table~\ref{tab:metadata}). Importantly, HLA haplotypes and clinical variables are \textbf{intentionally excluded} from the main classification pipeline to ensure the classifier relies only on learned gene-expression embeddings, not clinical labels. These variables are reserved for downstream biological interpretation.

\begin{table}[h!]
\centering
\caption{Clinical Metadata Columns in \texttt{seurat\_metadata.csv}}
\label{tab:metadata}
\begin{tabular}{ll}
\toprule
\textbf{Column} & \textbf{Description} \\
\midrule
\texttt{COND} & Disease label: \texttt{T1D} or \texttt{H} (Healthy) \\
\texttt{Sample\_ID} & Unique patient identifier \\
\texttt{Cluster\_Annotation\_Merged} & Immune cell type annotation \\
\texttt{Gender} & Patient sex \\
\texttt{Age\_at\_diagnosis} & Age at first T1D diagnosis \\
\texttt{Disease\_Duration} & Years since diagnosis \\
\texttt{DQ\_Risk\_Haplotypes} & HLA-DQ haplotype risk class \\
\bottomrule
\end{tabular}
\end{table}

\section{Cell Type Composition}

The annotated cell types span all major PBMC populations:

\begin{itemize}
    \item \textbf{T cells}: CD4+, CD8+, Regulatory T (Treg), Gamma-delta T
    \item \textbf{B cells}: Naive B, Memory B, Plasmablasts
    \item \textbf{NK cells}: Cytotoxic NK, NK-dim
    \item \textbf{Monocytes}: Classical (CD14+), Non-classical (CD16+)
    \item \textbf{Dendritic cells}: Plasmacytoid DC (pDC), Conventional DC (cDC)
\end{itemize}

\section{Celiac Disease Datasets}

To facilitate cross-disease comparison and validate the machine learning approach on a localised autoimmune condition, two Celiac disease datasets were included:

\subsection{Celiac scRNA-seq (GSE315138)}
Used for cross-disease latent space alignment. This dataset comprises single-cell RNA-seq data from PBMCs of 4 Celiac disease patients. The raw count matrix, featuring approximately 33,000 cells, was preprocessed using standard Seurat workflows (QC filtering, SCTransform normalisation) before projection into the T1D scVI latent space.

\subsection{Celiac Bulk RNA-seq (GSE145358)}
Used for robust machine learning classification. This dataset contains bulk transcriptomic profiles from 36 intestinal biopsy samples. The cohort includes three distinct patient states:
\begin{itemize}
    \item \textbf{Healthy Controls}: Normal intestinal mucosa
    \item \textbf{Silent Celiac (GFD)}: Patients adhering to a Gluten-Free Diet with recovered mucosa
    \item \textbf{Active Celiac}: Untreated patients with active villous atrophy
\end{itemize}
